\section{Kết luận}

Báo cáo này tập trung vào bước nền cho bài toán phân đoạn nguyên liệu thực phẩm fine-grained trên FoodSeg103 bằng cách chuẩn hóa baseline BiSeNet trực tiếp trên không gian nhãn nguyên thủy 103 lớp. Kết quả thực nghiệm cho thấy thiết lập này đủ để tạo một mốc tham chiếu rõ ràng, đồng thời làm lộ ra các hạn chế đặc trưng của mô hình lightweight khi xử lý vùng biên yếu, vùng chồng lấp và các lớp có texture gần nhau. Bên cạnh đó, nhánh thực nghiệm tái cấu trúc nhãn (103 xuống 75) được dùng như phân tích bổ sung để lượng hóa ảnh hưởng của phân mảnh nhãn, thay vì thay thế baseline chính. Trên cơ sở đó, bước tiếp theo của nghiên cứu là kiểm tra có hệ thống xem các mô-đun nhạy với texture, local detail hoặc boundary-aware có thể cải thiện chất lượng phân đoạn đến mức nào trên cùng thiết lập thực nghiệm.